\section{Processing variants and the robustness qualifier} \label{app:variants}

\subsection{Variants}

After the primary run, fixed alternative processings (variants) are run on the same product and reported next to the
primary result. They use the expected depths of the primary run; they are not recalibrated per variant. The variants
are:
\begin{enumerate}[label=(\alph*)]
\item detrending window 0.5\,d;
\item detrending window 1.0\,d;
\item cadences within $\pm 1$\,h of SPOC desaturation flags removed (only when the primary product is a SPOC product);
\item the QLP light curve with mask 1610619071: the first QLP release that is public when the other variants are run,
  or, if none is public then, the release that is public on 2027 February 28. It is run once and not re-run on later
  releases. It is not run when the primary product is itself QLP;
\item the \code{SAP\_FLUX} of the same product. This flux is not corrected for light from other stars, so variant (e)
  is diluted and leans toward ``shallower'';
\item for comparison only, outside the robustness rule: the counting rule of draft v2 (80\,\% coverage only, no edge
  zones, so every transit with 80\,\% coverage is counted).
\end{enumerate}

Variants are never edited. A variant that crashes or cannot be run is ``not available'' and is ignored. ``Not
available'' means one of two things: the registered command of the variant, run unchanged on the registered file,
stopped with an error while reading or computing (after the one repeat allowed by Appendix~\ref{app:procedure}, point
4); or its input file fails the time-system check of Appendix~\ref{app:p3}. Its attempt record
(Appendix~\ref{app:p4}), which contains the unedited error output or the printed header line, is saved and posted. A
refusal by the environment check or the file-name check is not ``not available'': the environment or the file name is
corrected and the variant is run. For variant (d), ``not yet public'' is handled by the robustness rule below.
\code{s110\_summary.py} copies the error output but cannot judge it. Because data and code are public, anyone can rerun
a variant declared not available. Appendix~\ref{app:p4} fixes what the saved record of an attempt must contain and when
a ``not available'' is invalid.

\subsection{Robustness qualifier} \label{app:robust}

Variants never change any primary outcome. They only decide whether an O1 or an O3 carries the qualifier ``robust''.
Only variants (a) to (e) with $\Sens \geq 0.95$ (qualifying variants) are considered.
\begin{itemize}
\item If no variant among (a) to (e) reaches $\Sens \geq 0.95$, the qualifier is ``robustness not testable''.
\item If at least one variant reaches $\Sens \geq 0.95$: a primary O1 is called robust only if no qualifying variant returns
  O3, and a primary O3 is called robust only if no qualifying variant returns O1 or O2.
\item O2 and O4 carry no robustness qualifier.
\end{itemize}

If variant (d) is not yet public when the other variants are run, ``robust'' and ``robustness not testable'' are
reported as provisional. A qualifying contradiction already gives ``not robust'', without a provisional qualifier:
a later variant cannot remove that contradiction. A provisional qualifier becomes final on 2027 February 28, with
(d) if (d) is public by then. ``Not yet public'' keeps the qualifier provisional; it
never makes it final. A provisional qualifier becomes final only in one of three cases: (d) has been run; a public QLP release
was attempted and the run failed; or a MAST search made on or after 2027 February 28 finds no QLP Sector 110 release.
\code{s110\_summary.py} accepts the last case only on or after that date (audit 6).

The qualifier is computed by \code{s110\_summary.py}, not by hand. The only exception is the cases of
Appendix~\ref{app:p5}, where the summary script cannot run; there the qualifier is written from the posted variant
results by exactly this rule. ``Robust'' has exactly the meaning given here, which is narrower than ``all variants
agree''. Appendix~\ref{app:p7} fixes this meaning and the line of variant outcomes that is always posted with it.

\subsection{Exploratory analyses}

Aperture photometry of my own and any search away from the predicted times (shifted transits, other periods) are
exploratory and cannot be cited for or against any outcome.
